\documentclass[../../main.tex]{subfiles}% 注意这里的文件路径不能用 ./main.tex ,否则用latexmk编译子文件会报错
\graphicspath{{\subfix{./image/}}{\subfix{../../image/}}} % 指定图片引用路径,后续可以直接使用图片文件名
% 如果图片存放在上级目录,则使用 ../../image/ 作为备用路径

% 例如:
% \begin{figure}[H]
% \centering
% \includegraphics[width=0.3\textwidth]{图.png}
% \caption{}
% \label{figure:图}
% \end{figure}
% 注意:上述\label{}一定要放在\caption{}之后,否则引用图片序号会只会显示??.

\begin{document}

\section{习题}

\begin{example}
利用推导 Laplace 方程的思想推导极小曲面方程。
\end{example}

\begin{example}
构造 $\mathbb{R}^n$ 上所有二次调和多项式组成的线性空间。
\end{example}

\begin{example}
仿照平均值公式的推导证明：当 $n\geqslant 3$ 时，对于第一边值问题
\begin{align*}
\begin{cases}
-\Delta u = f, & x\in B(0,r), \\
u = g, & x\in\partial B(0,r)
\end{cases}
\end{align*}
在 $C^2(B(0,r))\cap C^1(\overline{B}(0,r))$ 上的解 $u(x)$，成立
\begin{align*}
u(0) = \fint_{\partial B(0,r)} g(x)\,\mathrm{d}S(x) + \frac{1}{n(n-2)\alpha(n)} \int_{B(0,r)} \left( \frac{1}{|x|^{n-2}} - \frac{1}{r^{n-2}} \right) f(x)\,\mathrm{d}x.
\end{align*}
\end{example}

\begin{example}
仿照平均值公式的推导证明：当 $n=2$ 时，对于第一边值问题
\begin{align*}
\begin{cases}
-\Delta u = f, & x\in B(0,r), \\
u = g, & x\in\partial B(0,r)=C(0,r)
\end{cases}
\end{align*}
在 $C^2(B(0,r))\cap C^1(\overline{B}(0,r))$ 上的解 $u(x)$，成立
\begin{align*}
u(0) = \frac{1}{2\pi r}\int_{C(0,r)} g(x)\,\mathrm{d}l + \frac{1}{2\pi}\int_{B(0,r)} (\ln r - \ln |x|) f(x)\,\mathrm{d}x,
\end{align*}
其中 $\mathrm{d}l$ 为弧长微分。此时 $C(0,r)$ 是环绕圆盘 $B(0,r)$ 的圆周。
\end{example}

\begin{example}
若 $v\in C^2(\Omega)$ 满足
\begin{align*}
-\Delta v \leqslant 0,\quad x\in\Omega,
\end{align*}
则称 $v$ 在 $\Omega$ 上是下调和的。
\begin{enumerate}[(1)]
\item 证明：对于任意球 $B(x,r)\subset\Omega$，成立
\begin{align*}
v(x) \leqslant \fint_{B(x,r)} v(y)\,\mathrm{d}y.
\end{align*}
\item 证明：$\max_{\Omega} v(x) = \max_{\partial\Omega} v(x)$。
\item 设 $\phi:\mathbb{R}\to\mathbb{R}$ 是光滑凸函数，且 $u$ 是 $\Omega$ 上的调和函数。证明：$v=\phi(u)$ 是 $\Omega$ 上的下调和函数。
\item 设 $u$ 是 $\Omega$ 上的调和函数。证明：$v=|Du|^2$ 是 $\Omega$ 上的下调和函数。
\end{enumerate}
\end{example}

\begin{theorem}[Harnack定理]\label{theorem:Harnack定理}
假设 $\{u_n\}\subset C(\overline{\Omega})\cap C^2(\Omega)$ 是 $\Omega$ 上的调和函数列。如果 $\{u_n\}$ 在 $\partial\Omega$ 上一致收敛，则 $\{u_n\}$ 在 $\overline{\Omega}$ 上一致收敛，且收敛于一个调和函数。
\end{theorem}
\begin{note}
提示：利用极值原理和平均值公式.
\end{note}

\begin{theorem}[Schwarz 反射定理]\label{theorem:Schwarz 反射定理}
记上半球
\begin{align*}
B^+ = \{x=(x_1,x_2,\cdots,x_n)\in B(0,1) \mid x_n>0\},
\end{align*}
假设 $u$ 是上半球 $B^+$ 上的调和函数且在边界 $\{x\in\partial B^+ \mid x_n=0\}$ 上满足 $u=0$。令
\begin{align*}
v(x) = \begin{cases}
u(x_1,x_2,\cdots,x_n), & x_n\geqslant 0, \\
-u(x_1,x_2,\cdots,-x_n), & x_n<0.
\end{cases}
\end{align*}
证明：$v$ 是球 $B(0,1)$ 上的调和函数。
\end{theorem}
\begin{note}
提示：利用平均值公式的推论.
\end{note}

\begin{example}
设 $\Omega$ 是 $\mathbb{R}^n$ 的一个开集。如果原点 $0\notin\Omega$，我们记 $x^*=\frac{x}{|x|^2}$，$\Omega^*=\{x^* \mid x\in\Omega\}$。假设 $u$ 是在 $\Omega$ 上的一个函数，并定义 $\Omega^*$ 上的函数 $K[u](x)=|x|^{2-n}u(x^*)$，$x\in\Omega^*$ 为函数 $u$ 的 Kelvin 变换。证明：$u(x)$ 是 $\Omega$ 上的调和函数当且仅当 $K[u]$ 是 $\Omega^*$ 上的调和函数。
\end{example}

\begin{example}
设 $u(x)$ 是球 $B(0,R)$ 上的调和函数，且在 $\overline{B}(0,R)$ 上连续。又设
\begin{align*}
M = \int_{B(0,R)} u^2(x)\,\mathrm{d}x.
\end{align*}
试证：
\begin{enumerate}[(1)]
\item $\displaystyle |u(0)| \leqslant \left[ \frac{M}{\alpha(n)R^n} \right]^{\frac{1}{2}}$；
\item $\displaystyle |u(x)| \leqslant \left[ \frac{M}{\alpha(n)(R-|x|)^n} \right]^{\frac{1}{2}}$.
\end{enumerate}
\end{example}

\begin{example}
设 $u(x)$ 是单位球 $B=B(0,1)$ 上的有界调和函数。证明：
\begin{align*}
\sup_{x\in B} (1-|x|)|Du(x)| < +\infty.
\end{align*}
\end{example}
\begin{note}
提示：利用定理 \eqref{eq:der-est} 的第一步证明.
\end{note}

\begin{example}
设 $u(x)$ 是球 $B(0,R_0)$ 上的调和函数，对于 $R\in(0,R_0]$ 记
\begin{align*}
\omega(R) = \sup_{B(0,R)} u - \inf_{B(0,R)} u.
\end{align*}
\begin{enumerate}[(1)]
\item 利用 Harnack 不等式证明：存在 $\eta\in(0,1)$，使得
\begin{align*}
\omega\left(\frac{R}{2}\right) \leqslant \eta\,\omega(R).
\end{align*}
（提示：对调和函数 $w(x)=u(x)-\inf_{B(0,R)}u$ 在球 $B(0,\frac{R}{2})$ 上利用 Harnack 不等式）
\item 如果 $\sup_{B(0,R_0)} |u(x)| \leqslant M_0$，则存在常数 $\alpha\in(0,1)$，$C>0$，使得
\begin{align*}
\omega(R) \leqslant C(M_0+1)\left(\frac{R}{R_0}\right)^\alpha,\quad R\in(0,R_0].
\end{align*}
\end{enumerate}
\end{example}
\begin{note}
提示：对任意 $R\in(0,R_0)$，一定存在一个正整数 $i\geqslant 1$，使得 $\frac{R_0}{2^i}\leqslant R < \frac{R_0}{2^{i-1}}$.
\end{note}

\begin{theorem}[推广的Liouville定理]\label{theorem:推广的Liouville定理}
假设 $u$ 是 $\mathbb{R}^n$ 上的调和函数，且
\begin{align*}
|u(x)| \leqslant C_1|x|^m + C_2,\quad x\in\mathbb{R}^n,
\end{align*}
其中 $m$ 是非负整数，$C_1,C_2$ 是非负常数，则 $u$ 必为一个次数至多为 $m$ 的调和多项式。
\end{theorem}
\begin{note}
提示：只要证明 $D^{m+1}u(x)\equiv 0$.
\end{note}

\begin{example}
假设 $u\in C^2(\mathbb{R}^n)$。对于 $r>0$，定义
\begin{align*}
u_r(x) = \frac{1}{n\alpha(n)r^{n-1}} \int_{\partial B_r(x)} u(y)\,\mathrm{d}S(y).
\end{align*}
证明：
\begin{align*}
\Delta u_r = (\Delta u)_r.
\end{align*}
\end{example}

\begin{example}
证明：对于任意 $x=(x_1,x_2,\cdots,x_n)\in\mathbb{R}_+^n$，Poisson 核满足
\begin{align*}
\int_{\mathbb{R}^{n-1}} K(x,y)\,\mathrm{d}y = 1,
\end{align*}
其中
\begin{align*}
y=(y_1,y_2,\cdots,y_{n-1},0)\in\partial\mathbb{R}_+^n=\mathbb{R}^{n-1},\quad \mathrm{d}y=\mathrm{d}y_1\mathrm{d}y_2\cdots\mathrm{d}y_{n-1},\quad K(x,y)=\frac{2x_n}{n\alpha(n)|y-x|^n}.
\end{align*}
\end{example}

\begin{example}
求边值问题
\begin{align*}
\begin{cases}
-\Delta u = f, & x\in\Omega, \\
u|_{\partial\Omega} = g
\end{cases}
\end{align*}
的 Green 函数，其中
\begin{enumerate}[(1)]
\item $\Omega$ 是上半平面；
\item $\Omega$ 是第一象限；
\item $\Omega$ 是带形区域 $\{(x,y)\in\mathbb{R}^2 \mid x\in\mathbb{R},\ 0<y<l\}$，其中 $l$ 为正常数。
\end{enumerate}
\end{example}

\begin{example}
记 $B^+(R)=\{x=(x_1,x_2,\cdots,x_n)\in\mathbb{R}^n \mid x_n>0,\ |x|<R\}$ $(n\geqslant 2)$。求边值问题
\begin{align*}
\begin{cases}
-\Delta u = f, & x\in B^+(R), \\
u|_{\partial B^+(R)\cap\{x_n=0\}} = g, \\
u|_{\partial B^+(R)\cap\{|x|=R\}} = h
\end{cases}
\end{align*}
的 Green 函数。
\end{example}
\begin{note}
提示：分别考虑 $n=2$ 和 $n\geqslant 3$ 两种情形.
\end{note}

\begin{example}
证明：第二边值问题
\begin{align*}
\begin{cases}
u_{xx}+u_{yy}=0, & (x,y)\in B(0,R), \\
\left.\frac{\partial u}{\partial r}\right|_{r=R} = g(\theta), & \theta\in[0,2\pi]
\end{cases}
\end{align*}
的解在边值 $g(\theta)$ 满足条件 $\int_0^{2\pi}g(\theta)\,\mathrm{d}\theta=0$ 时可以表示成
\begin{align*}
u(r,\theta) = -\frac{1}{2\pi}\int_0^{2\pi} g(\tau) \ln\left[R^2+r^2-2Rr\cos(\tau-\theta)\right]\,\mathrm{d}\tau + C,
\end{align*}
其中 $C$ 为任意常数。这里 $(r,\theta)$ 是点 $(x,y)$ 的极坐标。
\end{example}

\begin{example}
假设 $u(x)\in C(\overline{\Omega})\cap C^2(\Omega)$ 是定解问题
\begin{align*}
\begin{cases}
-\Delta u + c(x)u = f(x), & x\in\Omega, \\
u|_{\partial\Omega} = 0
\end{cases}
\end{align*}
的一个解。
\begin{enumerate}[(1)]
\item 如果 $c(x)\geqslant c_0>0$，则
\begin{align*}
\max_{\overline{\Omega}} |u(x)| \leqslant \frac{1}{c_0}\sup_{\Omega}|f(x)|;
\end{align*}
\item 如果 $c(x)\geqslant 0$，则
\begin{align*}
\max_{\overline{\Omega}} |u(x)| \leqslant M\sup_{\Omega}|f(x)|,
\end{align*}
其中 $M$ 依赖于 $\Omega$ 的直径 $d$；（提示：不妨设原点 $0\in\Omega$，并令 $u(x)=w(x)v(x)=(d^2-|x|^2+1)v(x)$，然后考虑 $v(x)$ 满足的方程，利用 (1) 的证明方法可得 $v(x)$ 的最大模估计，从而得到 $u(x)$ 的最大模估计）
\item 如果 $c(x)<0$，试举反例说明上述最大模估计一般不成立。
\end{enumerate}
\end{example}

\begin{example}
假设 $u(x)\in C(\overline{\Omega})\cap C^2(\Omega)$ 是定解问题
\begin{align*}
\begin{cases}
-\Delta u = 1, & x\in\Omega, \\
u|_{\partial\Omega} = 0
\end{cases}
\end{align*}
的一个解。试证明：对于任意 $x_0\in\Omega$，估计
\begin{align*}
\frac{1}{2n}\min_{x\in\partial\Omega}|x-x_0|^2 \leqslant u(x_0) \leqslant \frac{1}{2n}\max_{x\in\partial\Omega}|x-x_0|^2
\end{align*}
成立，这里 $n$ 是空间的维数。
\end{example}

\begin{example}
假设 $u(x)\in C^1(\overline{\Omega})\cap C^2(\Omega)$ 是定解问题
\begin{align*}
\begin{cases}
-\Delta u + c(x)u = f(x), & x\in\Omega, \\
\left[\frac{\partial u}{\partial n}+\alpha(x)u\right]\bigg|_{\Gamma_1} = g_1, & u|_{\Gamma_2}=g_2
\end{cases}
\end{align*}
的一个解，其中 $\Gamma_1\cup\Gamma_2=\partial\Omega$，$\Gamma_1\cap\Gamma_2=\varnothing$，$\Gamma_2\neq\varnothing$。
\begin{enumerate}[(1)]
\item 如果 $c(x)\geqslant c_0>0$，$\alpha(x)\geqslant\alpha_0>0$，则有估计
\begin{align*}
\max_{\overline{\Omega}} |u(x)| \leqslant \max\left\{ \frac{1}{c_0}\sup_{\Omega}|f(x)|,\ \frac{1}{\alpha_0}\max_{\Gamma_1}|g_1|,\ \max_{\Gamma_2}|g_2| \right\}.
\end{align*}
\item 如果 $c(x)\geqslant 0$，$\alpha(x)\geqslant 0$，且 $\Gamma_1$ 满足内球条件，则上述问题的解是惟一的。这里 $\Gamma_1$ 满足内球条件是指对于任意 $x_0\in\Gamma_1$，存在一个球 $B$ 使得 $B\subset\Omega$，$\Gamma_1\cap\partial B=\{x_0\}$。（提示：令 $u(x)=w(x)v(x)$，其中 $w(x)$ 是待定的辅助函数。在球 $B$ 上我们考虑 $v(x)$ 满足的问题）
\end{enumerate}
\end{example}

\begin{example}
试用辅助函数
\begin{align*}
w(x)=e^{-a|x|^2}-e^{-aR^2}
\end{align*}
证明 Hopf 引理，这里 $a>0$，$R$ 是球 $B_R$ 的半径。
\end{example}

\begin{example}
假设 $\Omega\subset\mathbb{R}^n$ 是一个有界开集，$u_i(x)\in C^2(\Omega)\cap C(\overline{\Omega})$ $(i=1,2)$ 满足定解问题
\begin{align*}
\begin{cases}
-\Delta u_i + c_i(x)u_i = 0, & x\in\Omega, \\
u_i = g_i, & x\in\partial\Omega.
\end{cases}
\end{align*}
如果 $c_2(x)\geqslant c_1(x)\geqslant 0$，$g_1(x)\geqslant g_2(x)\geqslant 0$，则
\begin{align*}
u_1(x)\geqslant u_2(x).
\end{align*}
（提示：先证明 $u_i(x)\geqslant 0$，然后考虑 $w(x)=u_1(x)-u_2(x)$ 满足的定解问题）
\end{example}

\begin{example}
假设 $\Omega_0\subset\mathbb{R}^n$ 是一个有界区域，$\Omega=\mathbb{R}^n\setminus\overline{\Omega}_0$。如果 $u(x)\in C^2(\Omega)\cap C(\overline{\Omega})$ 是外部问题
\begin{align*}
\begin{cases}
-\Delta u + c(x)u = 0, & x\in\Omega, \\
u|_{\partial\Omega} = g(x), & x\in\partial\Omega, \\
\lim_{|x|\to\infty} u(x) = l
\end{cases}
\end{align*}
的一个解，其中 $c(x)\geqslant 0$ 且在 $\overline{\Omega}$ 上局部有界，则
\begin{align*}
\sup_{\Omega} |u(x)| \leqslant \max\{|l|,\ \max_{\partial\Omega}|g(x)|\}.
\end{align*}
（提示：先在区域 $B(0,R)\setminus\Omega_0$ 上证明关于 $u(x)$ 的估计，然后令 $R\to+\infty$）
\end{example}

\begin{example}
假设 $\Omega\subset\mathbb{R}^n$ 是一个有界开集。如果 $u(x)\in C^2(\Omega)\cap C(\overline{\Omega})$ 是定解问题
\begin{align*}
\begin{cases}
-\Delta u + |u|u = f, & x\in\Omega, \\
u|_{\partial\Omega} = g
\end{cases}
\end{align*}
的一个解，则
\begin{align*}
\max_{\overline{\Omega}} |u(x)| \leqslant \max\left\{ \max_{\partial\Omega}|g(x)|,\ \sup_{\Omega}|f|^{\frac{1}{2}} \right\}.
\end{align*}
\end{example}

\begin{example}
假设 $\Omega\subset\mathbb{R}^n$ 是一个有界开集，$u(x)\in C^2(\Omega)\cap C(\overline{\Omega})$ 是定解问题
\begin{align*}
\begin{cases}
-\Delta u + u^3 - u = 0, & x\in\Omega, \\
u|_{\partial\Omega} = g
\end{cases}
\end{align*}
的一个解。证明：如果 $\max_{\partial\Omega}|g(x)|\leqslant 1$，则 $\max_{\overline{\Omega}}|u(x)|\leqslant 1$。
\end{example}

\begin{example}
假设 $\Omega\subset\mathbb{R}^n$ 是一个有界开集，$u_i(x)\in C^2(\Omega)\cap C(\overline{\Omega})$ $(i=1,2)$ 满足定解问题
\begin{align*}
\begin{cases}
-\Delta u_i + u_i^3 = f_i(x), & x\in\Omega, \\
u_i = g_i(x), & x\in\partial\Omega.
\end{cases}
\end{align*}
如果 $f_1(x)\geqslant f_2(x)$，$g_1(x)\geqslant g_2(x)$，则
\begin{align*}
u_1(x)\geqslant u_2(x).
\end{align*}
\end{example}
\begin{note}
提示：考虑 $w(x)=u_1(x)-u_2(x)$ 满足的定解问题.
\end{note}

\begin{example}
假设 $\Omega\subset\mathbb{R}^n$ 是一个有界开集，$u(x)\in C^2(\Omega)\cap C(\overline{\Omega})$ 满足定解问题
\begin{align*}
\begin{cases}
-\Delta u + A\cdot Du = f(x), & x\in\Omega, \\
u = g(x), & x\in\partial\Omega,
\end{cases}
\end{align*}
其中 $A=A(x):\Omega\to\mathbb{R}^n$ 是一个有界连续向量函数。如果 $f(x)\geqslant 0$，$g(x)\geqslant 0$，则
\begin{align*}
u(x)\geqslant 0,\quad x\in\overline{\Omega}.
\end{align*}
（提示：不妨设原点 $0\in\Omega$。考虑 $w(x)=u(x)+\varepsilon(e^{Md}-e^{Mx_1})$ 满足的定解问题，其中 $M=\sup_{x\in\Omega}|A(x)|+1$，$d$ 为 $\Omega$ 的直径）
\end{example}

\begin{example}
假设 $\Omega\subset\mathbb{R}^n$ 是一个有界开集，$u_i(x)\in C^2(\Omega)\cap C^1(\overline{\Omega})$ $(i=1,2)$ 满足定解问题
\begin{align*}
\begin{cases}
-\Delta u_i + |Du_i|^2 = f_i(x), & x\in\Omega, \\
u_i = g_i(x), & x\in\partial\Omega.
\end{cases}
\end{align*}
如果 $f_1(x)\geqslant f_2(x)$，$g_1(x)\geqslant g_2(x)$，则
\begin{align*}
u_1(x)\geqslant u_2(x),\quad x\in\overline{\Omega}.
\end{align*}
（提示：考虑 $w(x)=u_1(x)-u_2(x)$ 满足的定解问题）
\end{example}

\begin{example}
假设 $\Omega\subset\mathbb{R}^n$ 是一个有界开集。如果 $u(x)\in C^2(\Omega)\cap C^1(\overline{\Omega})$ 是定解问题
\begin{align*}
\begin{cases}
-\Delta u + |u|^r u = f, & x\in\Omega, \\
\left[\frac{\partial u}{\partial n} + \alpha(x)u\right]\bigg|_{\partial\Omega} = g
\end{cases}
\end{align*}
的一个解，其中 $r>0$ 和 $\alpha(x)\geqslant\alpha_0>0$，则
\begin{align*}
\max_{\overline{\Omega}} |u(x)| \leqslant \max\left\{ \left(\sup_{\Omega}|f|\right)^{\frac{1}{r+1}},\ \frac{1}{\alpha_0}\max_{\partial\Omega}|g| \right\}.
\end{align*}
\end{example}

\begin{example}
假设 $\Omega\subset\mathbb{R}^n$ 是一个有界开集。如果 $u(x),v(x)\in C^2(\Omega)\cap C(\overline{\Omega})$ 满足方程组
\begin{align*}
\begin{cases}
-\Delta u + 2u - v = f, & x\in\Omega, \\
-\Delta v + 2v - u = g, & x\in\Omega
\end{cases}
\end{align*}
和边界条件
\begin{align*}
u|_{\partial\Omega} = v|_{\partial\Omega} = 0.
\end{align*}
证明：
\begin{align*}
\max\{ \max_{\Omega}|u(x)|,\ \max_{\Omega}|v(x)| \} \leqslant \max\{ \sup_{\Omega}|f(x)|,\ \sup_{\Omega}|g(x)| \}.
\end{align*}
\end{example}

\begin{example}
假设 $\Omega\subset\mathbb{R}^n$ 是一个有界开集，$x_0\in\partial\Omega$。如果 $u(x)\in C^2(\Omega)\cap C(\overline{\Omega}\setminus\{x_0\})$ 满足
\begin{align*}
\begin{cases}
-\Delta u = 0, & x\in\Omega, \\
u|_{\partial\Omega\setminus\{x_0\}} = g, \\
\lim_{x\to x_0} |u(x)| \leqslant M_0,
\end{cases}
\end{align*}
则
\begin{align*}
\sup_{\Omega} |u(x)| \leqslant \max\{M_0,\ \sup_{\partial\Omega\setminus\{x_0\}}|g(x)|\}.
\end{align*}
\end{example}

\begin{example}
记 $B^+=\{(x,y) \mid |x|^2+y^2<1,\ x\in\mathbb{R}^{n-1},\ y>0\}$ 是 $\mathbb{R}^n$ 的上半球。假设 $u(x,y)\in C^2(B^+)\cap C(\overline{B}^+)$ 是定解问题
\begin{align*}
\begin{cases}
-\Delta_x u - y u_{yy} + c(x,y)u = f(x,y), & (x,y)\in B^+, \\
u|_{\partial B^+} = g
\end{cases}
\end{align*}
的一个解。
\begin{enumerate}[(1)]
\item 如果 $c(x,y)\geqslant c_0>0$，则有估计
\begin{align*}
\max_{\overline{B}^+} |u(x,y)| \leqslant \max\{c_0^{-1}\sup_{B^+}|f(x,y)|,\ \max_{\partial B^+}|g(x,y)|\};
\end{align*}
\item 如果 $c(x,y)\geqslant 0$，则
\begin{align*}
\max_{\overline{B}^+} |u(x,y)| \leqslant M[\sup_{B^+}|f(x,y)| + \max_{\partial B^+}|g(x,y)|],
\end{align*}
其中 $M>0$ 是一常数.
\end{enumerate}
\end{example}
\begin{note}
提示：令 $u(x,y)=w(x)v(x,y)$，其中 $w(x)$ 是待定的辅助函数。在上半球 $B^+$ 上考虑 $v(x,y)$ 满足的问题，利用 (1) 的结论.
\end{note}

\begin{example}
记 $\mathbb{R}_+^2=\{(x,y) \mid x\in\mathbb{R},\ y>0\}$。证明：定解问题
\begin{align*}
\begin{cases}
-\Delta u = f(x,y), & (x,y)\in\mathbb{R}_+^2, \\
u|_{y=0} = g(x), & x\in\mathbb{R}
\end{cases}
\end{align*}
属于 $C^2(\mathbb{R}_+^2)\cap C(\overline{\mathbb{R}}_+^2)$ 的有界解是惟一的。（提示：考虑辅助函数 $w(x,y)=\varepsilon\ln[x^2+(y+1)^2]\pm u(x,y)$，其中 $\varepsilon$ 是任意正常数）
\end{example}

\begin{example}
假设 $\Omega\subset\mathbb{R}^n$ $(n\geqslant 3)$ 是一个有界开集，$x_0\in\partial\Omega$。如果 $u(x)\in C^2(\Omega)\cap C(\overline{\Omega}\setminus\{x_0\})$ 是定解问题
\begin{align*}
\begin{cases}
-\Delta u = f, & x\in\Omega, \\
u|_{\partial\Omega\setminus\{x_0\}} = g
\end{cases}
\end{align*}
的有界解。证明这样的解是惟一的。
\end{example}
\begin{note}
提示：考虑辅助函数 $w(x)=\frac{\varepsilon}{|x-x_0|^{n-2}}\pm u(x)$，其中 $d$ 是 $\Omega$ 的直径，$\varepsilon$ 是任意正常数.
\end{note}

\begin{example}
假设 $\Omega\subset\mathbb{R}^2$ 是一个有界区域，$x_0\in\partial\Omega$。如果 $u(x)\in C^2(\Omega)\cap C(\overline{\Omega}\setminus\{x_0\})$ 是定解问题
\begin{align*}
\begin{cases}
-\Delta u = f, & x\in\Omega, \\
u|_{\partial\Omega\setminus\{x_0\}} = g
\end{cases}
\end{align*}
的有界解。证明这样的解是惟一的。
\end{example}
\begin{note}
提示：考虑辅助函数 $w(x)=\varepsilon\ln\left(\frac{d}{|x-x_0|}\right)\pm u(x)$，其中 $d$ 是 $\Omega$ 的直径，$\varepsilon$ 是任意正常数.
\end{note}

\begin{example}
假设 $\Omega\subset\mathbb{R}^n$ $(n\geqslant 3)$ 是一个有界开集，$x_0\in\Omega$。又设 $u(x)\in C^2(\Omega)\cap C(\overline{\Omega})$ 是定解问题
\begin{align*}
\begin{cases}
-\Delta u = f, & x\in\Omega, \\
u|_{\partial\Omega} = g
\end{cases}
\end{align*}
的解，$v(x)\in C^2(\Omega)\cap C(\overline{\Omega}\setminus\{x_0\})$ 是定解问题
\begin{align*}
\begin{cases}
-\Delta v = f, & x\in\Omega, \\
v|_{\partial\Omega} = g
\end{cases}
\end{align*}
的有界解。证明：$x_0$ 是 $v(x)$ 的可去奇点，即
\begin{align*}
u(x)\equiv v(x),\quad x\in\overline{\Omega}\setminus\{x_0\}.
\end{align*}
（提示：当 $n\geqslant 3$ 时，考虑辅助函数 $w(x)=\frac{\varepsilon}{|x-x_0|^{n-2}}\pm [u(x)-v(x)]$，其中 $\varepsilon$ 是任意正常数。当 $n=2$ 时，考虑辅助函数 $w(x)=-\varepsilon\ln|x-x_0|\pm [u(x)-v(x)]$）
\end{example}

\begin{example}
假设 $\Omega\subset\mathbb{R}^n$ $(n\geqslant 3)$ 是一个有界开集。考虑定解问题
\begin{align*}
\begin{cases}
-\Delta u(x) + A(x)\cdot Du(x) + c(x)u(x) = f(x), & x\in\Omega, \\
u|_{\partial\Omega} = 0,
\end{cases}
\end{align*}
其中 $n$ 维向量函数 $A:\Omega\to\mathbb{R}^n$ 和函数 $c(x)$ 在 $\Omega$ 上连续有界。如果条件 $c(x)-\frac{1}{4}|A(x)|^2>0$ 成立，利用能量估计方法证明上述问题解的唯一性。
\end{example}

\begin{example}
假设 $\Omega\subset\mathbb{R}^n$ $(n\geqslant 3)$ 是一个有界开集，且在 $\Omega$ 上 $c(x)\geqslant 0$，试利用能量估计方法证明 Neumann 边值问题
\begin{align*}
\begin{cases}
-\Delta u(x) + c(x)u(x) = f(x), & x\in\Omega, \\
\frac{\partial u}{\partial n} = g(x), & x\in\partial\Omega
\end{cases}
\end{align*}
在函数类 $C^2(\Omega)\cap C^1(\overline{\Omega})$ 中的解在相差一个常数的意义下惟一。
\end{example}



\end{document}